\documentclass[11pt, oneside]{book}
\usepackage[margin=1in]{geometry}
\usepackage{amsmath,amssymb,amsthm}
\usepackage{hyperref}
\usepackage{graphicx}
\usepackage{booktabs}
\usepackage{setspace}
\onehalfspacing

\usepackage[backend=biber,style=authoryear]{biblatex}
\addbibresource{superluminal_bibliography.bib}

\title{Superluminal Phenomena in Relativistic Astrophysics:\\
Modeling, Inference, and Simulation Workflows}
\author{Integrated Pass-2 to Pass-4 Research Pack}
\date{2026-02-23}

\begin{document}
\frontmatter
\maketitle
\tableofcontents

\chapter{Executive Synthesis}
Apparent superluminal motion in relativistic outflows is now treated as a geometric and radiative projection effect, not physical superluminal signaling. The central modeling constraint is therefore:
\begin{equation}
\beta_{\mathrm{app}} = \frac{\beta \sin\theta}{1-\beta\cos\theta}
\end{equation}
with \(\beta\le 1\), while \(\beta_{\mathrm{app}}\) can exceed unity for high \(\Gamma\) and small \(\theta\).  

This deliverable binds three implementation layers:
\begin{itemize}
  \item \textbf{Geometric observer layer}: transform intrinsic motions to apparent observables.
  \item \textbf{Channel transport layer}: electromagnetic/GW/CR/neutrino arrival differences, including delays.
  \item \textbf{Physics closure layer}: shock/jets/transport models and optional GRMHD synthesis.
\end{itemize}

\mainmatter
\chapter{Pass-2 Enriched Source Stack}
The enriched pass-2 catalog extends historical and modern references with:
\begin{enumerate}
  \item DOI/arXiv dual identifiers.
  \item Explicit modeling assumptions.
  \item Reproducibility notes (dataset + code entry points).
\end{enumerate}
Core references include Rees (1966), Blandford-Rees (1977), Blandford-K{\"o}nigl (1979), and Urry-Padovani (1995).  

Current best practice is to validate every claim by mapping it to:
\begin{itemize}
  \item a formal publication record,
  \item open data,
  \item executable inference setup.
\end{itemize}

\section{Observed anchors}
MOJAVE VI--XVIII provide the long-baseline kinematics baseline for apparent-speed distributions.  
GW170817 multi-band studies constrain causal timing and structured-jet interpretation.  
TXS--neutrino and UHECR source-correlation work remains hypothesis-weighted and should be carried with uncertainty tagging.

\chapter{Pass-3 Workflow and Implementation}
Pass-3 captures a practical computational design that is intentionally split by physics layer:

\section{1. Geometry-first baseline}
Event tuples are created as
\((t_0,\mathbf r_0,\beta,\Gamma,\theta,P(t),\phi_c)\), then transformed to observer frame by finite-speed delay and Doppler factors.

\section{2. Multi-channel transport}
Each channel uses explicit delay kernels \(K_c(t)\) and attenuation terms.  
Channel ordering is emergent:
\begin{itemize}
  \item GW: near-luminal causal promptness.
  \item EM: opacity + jet-structure delays.
  \item CR/neutrino: larger dispersion and transport uncertainty.
\end{itemize}

\section{3. Multi-framework execution}
\begin{itemize}
  \item \textbf{Python path}: `numpy`/`scipy` + `numpyro`/`emcee` for workflow and fitting.
  \item \textbf{Julia path}: `DifferentialEquations.jl` + `Turing.jl` for Bayesian/inference workflows.
  \item \textbf{MHD path}: HARM/PLUTO/BHAC post-processing with synthetic-emission mapping.
\end{itemize}

\section{4. Diagnostics and reproducibility}
Mandatory outputs include:
\begin{itemize}
  \item posterior corner plots,
  \item \(\beta_{\mathrm{app}}(t)\) residual trajectories,
  \item channel arrival lag distributions,
  \item run manifests with data hash and seed state.
\end{itemize}

\chapter{Mathematical Model Inventory}
The mathematical toolset includes:
\begin{enumerate}
  \item Light-travel and Doppler geometry.
  \item Pattern-speed versus bulk-flow decomposition.
  \item Radiative transfer with ray integration.
  \item CR/neutrino transport equation classes (advection-diffusion-loss source form).
  \item Bayesian objective fitting with structured priors.
  \item Optional speculative couplers (e.g., high-order GW-EM conversion probes in low-impact mode).
\end{enumerate}

Representative modeling equations (selection):
\[
\delta=[\Gamma(1-\beta\cos\theta)]^{-1},\quad
S_{\nu,\mathrm{obs}}=\delta^{k-\alpha}S'_{\nu'}
\]
\[
\frac{\partial f_c}{\partial t}+\mathbf v\cdot\nabla f_c=
\nabla\cdot(\mathbf D\nabla f_c)+\frac13(\nabla\cdot\mathbf u)\frac{\partial f_c}{\partial\ln p}+Q_c-L_c
\]
\[
\mathcal L(\theta)=\sum_{t,\nu}\frac{\left[X_{\mathrm{model}}(t,\nu;\theta)-X_{\mathrm{obs}}(t,\nu)\right]^2}{\sigma^2}
\]

\chapter{Reference Integrity and Citation}
For publication-grade writing, the preferred source-of-truth is:
\begin{itemize}
  \item `superluminal_bibliography.bib`
  \item DOI links for foundational references
  \item arXiv links where DOI is absent or behind archival barriers
\end{itemize}

\printbibliography
\end{document}

